\section{27.11.2025 -- Arbeitsrecht und Verträge}

\begin{vocabtable}
\deword{geltend}{valid, in force, applicable}{Die Regelung ist weiterhin geltend.}
\deword{gerechtfertigt}{justified, warranted}{Die Kündigung war nicht gerechtfertigt.}
\deword{unwirksam}{invalid, not legally effective}{Die Kündigung kann unwirksam sein.}
\deword{innerhalb}{within}{Innerhalb von drei Wochen muss man reagieren.}
\deword{der Zugang}{receipt, delivery of a document}{Der Zugang der Kündigung ist wichtig.}
\deword{die Klage}{lawsuit, legal action}{Er erhebt Klage beim Arbeitsgericht.}
\deword{das Arbeitsgericht}{labor court}{Das Arbeitsgericht entscheidet über die Klage.}
\deword{die Feststellung}{determination, declaration, finding}{Die Feststellung erfolgt durch das Gericht.}
\deword{erheben}{to file, to raise}{Man kann eine Klage erheben.}
\deword{das Arbeitsverhältnis}{employment relationship}{Das Arbeitsverhältnis wurde beendet.}
\deword{aufgelöst}{dissolved, terminated}{Der Vertrag wurde aufgelöst.}
\deword{die Überschrift}{heading, title}{Die Überschrift ist klar formuliert.}
\deword{überfliegen}{to skim}{Ich habe den Text zuerst überflogen.}
\deword{die Abschnitte}{sections, paragraphs}{Der Vertrag hat mehrere Abschnitte.}
\deword{die Probezeit}{probationary period}{Die Probezeit dauert sechs Monate.}
\deword{unbedingt}{absolutely, definitely}{Das muss man unbedingt beachten.}
\deword{des Arbeitsverhältnisses}{of the employment relationship}{Die Beendigung des Arbeitsverhältnisses ist geregelt.}
\deword{die Kündigungsfrist}{notice period}{Die Kündigungsfrist beträgt vier Wochen.}
\deword{die Kündigungsschutzklage}{unfair-dismissal claim}{Eine Kündigungsschutzklage muss schnell erhoben werden.}
\deword{die Nebenabreden}{side agreements}{Nebenabreden müssen schriftlich festgehalten werden.}
\deword{die Verschwiegenheitspflicht}{duty of confidentiality}{Die Verschwiegenheitspflicht gilt auch nach Vertragsende.}
\deword{verschwiegen}{discreet, tight-lipped}{Er ist sehr verschwiegen.}
\deword{verpflichtet}{obliged, required, bound}{Ich bin zur Geheimhaltung verpflichtet.}
\deword{entgeltlich}{paid, remunerated}{Eine entgeltliche Nebentätigkeit muss genehmigt werden.}
\deword{beeinträchtigend}{impairing, detrimental}{Die Nebentätigkeit darf die Arbeit nicht beeinträchtigen.}
\deword{die Zustimmung}{consent, approval}{Dafür braucht man die Zustimmung des Arbeitgebers.}
\deword{befolgen}{to comply with, to obey}{Man muss die Anweisungen befolgen.}
\deword{die Überraschung}{surprise}{Die Entscheidung war eine Überraschung.}
\deword{die Beobachtung}{observation, monitoring}{Die Beobachtung war interessant.}
\deword{zuständig}{responsible, in charge, competent}{Wer ist dafür zuständig?}
\deword{vertrauen}{to trust}{Ich vertraue meinem Team.}
\deword{die Bedingungen}{conditions, terms, requirements}{Die Bedingungen stehen im Vertrag.}
\deword{der Betriebsrat}{works council}{Der Betriebsrat vertritt die Beschäftigten.}
\deword{die Mitgliedschaft}{membership}{Die Mitgliedschaft ist freiwillig.}
\deword{die Abfindung}{severance payment}{Nach der Kündigung erhielt er eine Abfindung.}
\deword{entlassen}{to dismiss, to fire, to release}{Der Mitarbeiter wurde entlassen.}
\end{vocabtable}
